% F083｜数学提要；12式逐式TikZ图解见02_explanation.tex，图源归assets/tikz/F083_*。
\textbf{结论名称：}F083 常见指数、对数型函数的奇偶性速判

\textbf{解决的问题：}看到指数、对数、根号与分式交织的函数，不必一一硬算，就能判断是奇函数还是偶函数，并记住必须检查的定义域条件。

\textbf{第一道关：}必须保证定义域关于 $0$ 对称：如果 $x$ 能取，$-x$ 也必须能取。然后比较 $f(-x)$ 和 $f(x)$：
\[
\boxed{f(-x)=-f(x)\Rightarrow\text{奇函数}},\qquad
\boxed{f(-x)=f(x)\Rightarrow\text{偶函数}}.
\]

\textbf{第二道关：}找出变号后最简单的变化：
\[
a^{-x}=\frac1{a^x},\qquad
\sqrt{x^2+1}-x=\frac1{\sqrt{x^2+1}+x},\qquad
\ln\frac1t=-\ln t\ (t>0).
\]
对于 $a>0,a\ne1$，下列函数的定义域需分别判断。

\textbf{常见的 7 个奇函数：}
\[
\begin{gathered}
 a^x-a^{-x},\quad \frac{a^x-1}{a^x+1},\quad \frac{a^x+1}{a^x-1},\\
 \ln\frac{kx-b}{kx+b},\quad \ln\frac{kx+b}{kx-b},\\
 \log_a\!\bigl(x+\sqrt{x^2+1}\bigr),\quad
 \log_a\!\bigl(\sqrt{x^2+1}-x\bigr).
\end{gathered}
\]
其中两种指数分式的第三个函数须 $x\ne0$；两种对数分式在 $k,b\ne0$ 时须 $|x|>|b/k|$；两个根式对数在实数范围都有定义。

\textbf{常见的 5 个偶函数：}
\[
\begin{gathered}
 a^x+a^{-x},\quad a^{|x|},\quad a^{-x^2},\quad \log_a|x|,\\
 x\left(\frac12+\frac1{2^x-1}\right).
\end{gathered}
\]
最后两个函数在 $x=0$ 无定义；特别是最后一个，不能代入 $x=0$。

\textbf{一眼记法：}“指数成对，一加偶、一减奇；倒数互换，分式或对数会变号；绝对值或平方常让函数变偶；两个奇函数相乘是偶函数。”

\textbf{使用提醒：}这是一组\textbf{典型函数的速判规律}，不是看到指数或对数就一定有奇偶性。参数和定义域一旦变化，都要重新核对。
